\documentclass[10pt]{article} % For LaTeX2e
\usepackage{tmlr}
% If accepted, instead use the following line for the camera-ready submission:
%\usepackage[accepted]{tmlr}
% To de-anonymize and remove mentions to TMLR (for example for posting to preprint servers), instead use the following:
%\usepackage[preprint]{tmlr}

% Optional math commands from https://github.com/goodfeli/dlbook_notation.
%%%%% NEW MATH DEFINITIONS %%%%%

\usepackage{amsmath,amsfonts,bm}

% Mark sections of captions for referring to divisions of figures
\newcommand{\figleft}{{\em (Left)}}
\newcommand{\figcenter}{{\em (Center)}}
\newcommand{\figright}{{\em (Right)}}
\newcommand{\figtop}{{\em (Top)}}
\newcommand{\figbottom}{{\em (Bottom)}}
\newcommand{\captiona}{{\em (a)}}
\newcommand{\captionb}{{\em (b)}}
\newcommand{\captionc}{{\em (c)}}
\newcommand{\captiond}{{\em (d)}}

% Highlight a newly defined term
\newcommand{\newterm}[1]{{\bf #1}}


% Figure reference, lower-case.
\def\figref#1{figure~\ref{#1}}
% Figure reference, capital. For start of sentence
\def\Figref#1{Figure~\ref{#1}}
\def\twofigref#1#2{figures \ref{#1} and \ref{#2}}
\def\quadfigref#1#2#3#4{figures \ref{#1}, \ref{#2}, \ref{#3} and \ref{#4}}
% Section reference, lower-case.
\def\secref#1{section~\ref{#1}}
% Section reference, capital.
\def\Secref#1{Section~\ref{#1}}
% Reference to two sections.
\def\twosecrefs#1#2{sections \ref{#1} and \ref{#2}}
% Reference to three sections.
\def\secrefs#1#2#3{sections \ref{#1}, \ref{#2} and \ref{#3}}
% Reference to an equation, lower-case.
\def\eqref#1{equation~\ref{#1}}
% Reference to an equation, upper case
\def\Eqref#1{Equation~\ref{#1}}
% A raw reference to an equation---avoid using if possible
\def\plaineqref#1{\ref{#1}}
% Reference to a chapter, lower-case.
\def\chapref#1{chapter~\ref{#1}}
% Reference to an equation, upper case.
\def\Chapref#1{Chapter~\ref{#1}}
% Reference to a range of chapters
\def\rangechapref#1#2{chapters\ref{#1}--\ref{#2}}
% Reference to an algorithm, lower-case.
\def\algref#1{algorithm~\ref{#1}}
% Reference to an algorithm, upper case.
\def\Algref#1{Algorithm~\ref{#1}}
\def\twoalgref#1#2{algorithms \ref{#1} and \ref{#2}}
\def\Twoalgref#1#2{Algorithms \ref{#1} and \ref{#2}}
% Reference to a part, lower case
\def\partref#1{part~\ref{#1}}
% Reference to a part, upper case
\def\Partref#1{Part~\ref{#1}}
\def\twopartref#1#2{parts \ref{#1} and \ref{#2}}

\def\ceil#1{\lceil #1 \rceil}
\def\floor#1{\lfloor #1 \rfloor}
\def\1{\bm{1}}
\newcommand{\train}{\mathcal{D}}
\newcommand{\valid}{\mathcal{D_{\mathrm{valid}}}}
\newcommand{\test}{\mathcal{D_{\mathrm{test}}}}

\def\eps{{\epsilon}}


% Random variables
\def\reta{{\textnormal{$\eta$}}}
\def\ra{{\textnormal{a}}}
\def\rb{{\textnormal{b}}}
\def\rc{{\textnormal{c}}}
\def\rd{{\textnormal{d}}}
\def\re{{\textnormal{e}}}
\def\rf{{\textnormal{f}}}
\def\rg{{\textnormal{g}}}
\def\rh{{\textnormal{h}}}
\def\ri{{\textnormal{i}}}
\def\rj{{\textnormal{j}}}
\def\rk{{\textnormal{k}}}
\def\rl{{\textnormal{l}}}
% rm is already a command, just don't name any random variables m
\def\rn{{\textnormal{n}}}
\def\ro{{\textnormal{o}}}
\def\rp{{\textnormal{p}}}
\def\rq{{\textnormal{q}}}
\def\rr{{\textnormal{r}}}
\def\rs{{\textnormal{s}}}
\def\rt{{\textnormal{t}}}
\def\ru{{\textnormal{u}}}
\def\rv{{\textnormal{v}}}
\def\rw{{\textnormal{w}}}
\def\rx{{\textnormal{x}}}
\def\ry{{\textnormal{y}}}
\def\rz{{\textnormal{z}}}

% Random vectors
\def\rvepsilon{{\mathbf{\epsilon}}}
\def\rvtheta{{\mathbf{\theta}}}
\def\rva{{\mathbf{a}}}
\def\rvb{{\mathbf{b}}}
\def\rvc{{\mathbf{c}}}
\def\rvd{{\mathbf{d}}}
\def\rve{{\mathbf{e}}}
\def\rvf{{\mathbf{f}}}
\def\rvg{{\mathbf{g}}}
\def\rvh{{\mathbf{h}}}
\def\rvu{{\mathbf{i}}}
\def\rvj{{\mathbf{j}}}
\def\rvk{{\mathbf{k}}}
\def\rvl{{\mathbf{l}}}
\def\rvm{{\mathbf{m}}}
\def\rvn{{\mathbf{n}}}
\def\rvo{{\mathbf{o}}}
\def\rvp{{\mathbf{p}}}
\def\rvq{{\mathbf{q}}}
\def\rvr{{\mathbf{r}}}
\def\rvs{{\mathbf{s}}}
\def\rvt{{\mathbf{t}}}
\def\rvu{{\mathbf{u}}}
\def\rvv{{\mathbf{v}}}
\def\rvw{{\mathbf{w}}}
\def\rvx{{\mathbf{x}}}
\def\rvy{{\mathbf{y}}}
\def\rvz{{\mathbf{z}}}

% Elements of random vectors
\def\erva{{\textnormal{a}}}
\def\ervb{{\textnormal{b}}}
\def\ervc{{\textnormal{c}}}
\def\ervd{{\textnormal{d}}}
\def\erve{{\textnormal{e}}}
\def\ervf{{\textnormal{f}}}
\def\ervg{{\textnormal{g}}}
\def\ervh{{\textnormal{h}}}
\def\ervi{{\textnormal{i}}}
\def\ervj{{\textnormal{j}}}
\def\ervk{{\textnormal{k}}}
\def\ervl{{\textnormal{l}}}
\def\ervm{{\textnormal{m}}}
\def\ervn{{\textnormal{n}}}
\def\ervo{{\textnormal{o}}}
\def\ervp{{\textnormal{p}}}
\def\ervq{{\textnormal{q}}}
\def\ervr{{\textnormal{r}}}
\def\ervs{{\textnormal{s}}}
\def\ervt{{\textnormal{t}}}
\def\ervu{{\textnormal{u}}}
\def\ervv{{\textnormal{v}}}
\def\ervw{{\textnormal{w}}}
\def\ervx{{\textnormal{x}}}
\def\ervy{{\textnormal{y}}}
\def\ervz{{\textnormal{z}}}

% Random matrices
\def\rmA{{\mathbf{A}}}
\def\rmB{{\mathbf{B}}}
\def\rmC{{\mathbf{C}}}
\def\rmD{{\mathbf{D}}}
\def\rmE{{\mathbf{E}}}
\def\rmF{{\mathbf{F}}}
\def\rmG{{\mathbf{G}}}
\def\rmH{{\mathbf{H}}}
\def\rmI{{\mathbf{I}}}
\def\rmJ{{\mathbf{J}}}
\def\rmK{{\mathbf{K}}}
\def\rmL{{\mathbf{L}}}
\def\rmM{{\mathbf{M}}}
\def\rmN{{\mathbf{N}}}
\def\rmO{{\mathbf{O}}}
\def\rmP{{\mathbf{P}}}
\def\rmQ{{\mathbf{Q}}}
\def\rmR{{\mathbf{R}}}
\def\rmS{{\mathbf{S}}}
\def\rmT{{\mathbf{T}}}
\def\rmU{{\mathbf{U}}}
\def\rmV{{\mathbf{V}}}
\def\rmW{{\mathbf{W}}}
\def\rmX{{\mathbf{X}}}
\def\rmY{{\mathbf{Y}}}
\def\rmZ{{\mathbf{Z}}}

% Elements of random matrices
\def\ermA{{\textnormal{A}}}
\def\ermB{{\textnormal{B}}}
\def\ermC{{\textnormal{C}}}
\def\ermD{{\textnormal{D}}}
\def\ermE{{\textnormal{E}}}
\def\ermF{{\textnormal{F}}}
\def\ermG{{\textnormal{G}}}
\def\ermH{{\textnormal{H}}}
\def\ermI{{\textnormal{I}}}
\def\ermJ{{\textnormal{J}}}
\def\ermK{{\textnormal{K}}}
\def\ermL{{\textnormal{L}}}
\def\ermM{{\textnormal{M}}}
\def\ermN{{\textnormal{N}}}
\def\ermO{{\textnormal{O}}}
\def\ermP{{\textnormal{P}}}
\def\ermQ{{\textnormal{Q}}}
\def\ermR{{\textnormal{R}}}
\def\ermS{{\textnormal{S}}}
\def\ermT{{\textnormal{T}}}
\def\ermU{{\textnormal{U}}}
\def\ermV{{\textnormal{V}}}
\def\ermW{{\textnormal{W}}}
\def\ermX{{\textnormal{X}}}
\def\ermY{{\textnormal{Y}}}
\def\ermZ{{\textnormal{Z}}}

% Vectors
\def\vzero{{\bm{0}}}
\def\vone{{\bm{1}}}
\def\vmu{{\bm{\mu}}}
\def\vtheta{{\bm{\theta}}}
\def\va{{\bm{a}}}
\def\vb{{\bm{b}}}
\def\vc{{\bm{c}}}
\def\vd{{\bm{d}}}
\def\ve{{\bm{e}}}
\def\vf{{\bm{f}}}
\def\vg{{\bm{g}}}
\def\vh{{\bm{h}}}
\def\vi{{\bm{i}}}
\def\vj{{\bm{j}}}
\def\vk{{\bm{k}}}
\def\vl{{\bm{l}}}
\def\vm{{\bm{m}}}
\def\vn{{\bm{n}}}
\def\vo{{\bm{o}}}
\def\vp{{\bm{p}}}
\def\vq{{\bm{q}}}
\def\vr{{\bm{r}}}
\def\vs{{\bm{s}}}
\def\vt{{\bm{t}}}
\def\vu{{\bm{u}}}
\def\vv{{\bm{v}}}
\def\vw{{\bm{w}}}
\def\vx{{\bm{x}}}
\def\vy{{\bm{y}}}
\def\vz{{\bm{z}}}

% Elements of vectors
\def\evalpha{{\alpha}}
\def\evbeta{{\beta}}
\def\evepsilon{{\epsilon}}
\def\evlambda{{\lambda}}
\def\evomega{{\omega}}
\def\evmu{{\mu}}
\def\evpsi{{\psi}}
\def\evsigma{{\sigma}}
\def\evtheta{{\theta}}
\def\eva{{a}}
\def\evb{{b}}
\def\evc{{c}}
\def\evd{{d}}
\def\eve{{e}}
\def\evf{{f}}
\def\evg{{g}}
\def\evh{{h}}
\def\evi{{i}}
\def\evj{{j}}
\def\evk{{k}}
\def\evl{{l}}
\def\evm{{m}}
\def\evn{{n}}
\def\evo{{o}}
\def\evp{{p}}
\def\evq{{q}}
\def\evr{{r}}
\def\evs{{s}}
\def\evt{{t}}
\def\evu{{u}}
\def\evv{{v}}
\def\evw{{w}}
\def\evx{{x}}
\def\evy{{y}}
\def\evz{{z}}

% Matrix
\def\mA{{\bm{A}}}
\def\mB{{\bm{B}}}
\def\mC{{\bm{C}}}
\def\mD{{\bm{D}}}
\def\mE{{\bm{E}}}
\def\mF{{\bm{F}}}
\def\mG{{\bm{G}}}
\def\mH{{\bm{H}}}
\def\mI{{\bm{I}}}
\def\mJ{{\bm{J}}}
\def\mK{{\bm{K}}}
\def\mL{{\bm{L}}}
\def\mM{{\bm{M}}}
\def\mN{{\bm{N}}}
\def\mO{{\bm{O}}}
\def\mP{{\bm{P}}}
\def\mQ{{\bm{Q}}}
\def\mR{{\bm{R}}}
\def\mS{{\bm{S}}}
\def\mT{{\bm{T}}}
\def\mU{{\bm{U}}}
\def\mV{{\bm{V}}}
\def\mW{{\bm{W}}}
\def\mX{{\bm{X}}}
\def\mY{{\bm{Y}}}
\def\mZ{{\bm{Z}}}
\def\mBeta{{\bm{\beta}}}
\def\mPhi{{\bm{\Phi}}}
\def\mLambda{{\bm{\Lambda}}}
\def\mSigma{{\bm{\Sigma}}}

% Tensor
\DeclareMathAlphabet{\mathsfit}{\encodingdefault}{\sfdefault}{m}{sl}
\SetMathAlphabet{\mathsfit}{bold}{\encodingdefault}{\sfdefault}{bx}{n}
\newcommand{\tens}[1]{\bm{\mathsfit{#1}}}
\def\tA{{\tens{A}}}
\def\tB{{\tens{B}}}
\def\tC{{\tens{C}}}
\def\tD{{\tens{D}}}
\def\tE{{\tens{E}}}
\def\tF{{\tens{F}}}
\def\tG{{\tens{G}}}
\def\tH{{\tens{H}}}
\def\tI{{\tens{I}}}
\def\tJ{{\tens{J}}}
\def\tK{{\tens{K}}}
\def\tL{{\tens{L}}}
\def\tM{{\tens{M}}}
\def\tN{{\tens{N}}}
\def\tO{{\tens{O}}}
\def\tP{{\tens{P}}}
\def\tQ{{\tens{Q}}}
\def\tR{{\tens{R}}}
\def\tS{{\tens{S}}}
\def\tT{{\tens{T}}}
\def\tU{{\tens{U}}}
\def\tV{{\tens{V}}}
\def\tW{{\tens{W}}}
\def\tX{{\tens{X}}}
\def\tY{{\tens{Y}}}
\def\tZ{{\tens{Z}}}


% Graph
\def\gA{{\mathcal{A}}}
\def\gB{{\mathcal{B}}}
\def\gC{{\mathcal{C}}}
\def\gD{{\mathcal{D}}}
\def\gE{{\mathcal{E}}}
\def\gF{{\mathcal{F}}}
\def\gG{{\mathcal{G}}}
\def\gH{{\mathcal{H}}}
\def\gI{{\mathcal{I}}}
\def\gJ{{\mathcal{J}}}
\def\gK{{\mathcal{K}}}
\def\gL{{\mathcal{L}}}
\def\gM{{\mathcal{M}}}
\def\gN{{\mathcal{N}}}
\def\gO{{\mathcal{O}}}
\def\gP{{\mathcal{P}}}
\def\gQ{{\mathcal{Q}}}
\def\gR{{\mathcal{R}}}
\def\gS{{\mathcal{S}}}
\def\gT{{\mathcal{T}}}
\def\gU{{\mathcal{U}}}
\def\gV{{\mathcal{V}}}
\def\gW{{\mathcal{W}}}
\def\gX{{\mathcal{X}}}
\def\gY{{\mathcal{Y}}}
\def\gZ{{\mathcal{Z}}}

% Sets
\def\sA{{\mathbb{A}}}
\def\sB{{\mathbb{B}}}
\def\sC{{\mathbb{C}}}
\def\sD{{\mathbb{D}}}
% Don't use a set called E, because this would be the same as our symbol
% for expectation.
\def\sF{{\mathbb{F}}}
\def\sG{{\mathbb{G}}}
\def\sH{{\mathbb{H}}}
\def\sI{{\mathbb{I}}}
\def\sJ{{\mathbb{J}}}
\def\sK{{\mathbb{K}}}
\def\sL{{\mathbb{L}}}
\def\sM{{\mathbb{M}}}
\def\sN{{\mathbb{N}}}
\def\sO{{\mathbb{O}}}
\def\sP{{\mathbb{P}}}
\def\sQ{{\mathbb{Q}}}
\def\sR{{\mathbb{R}}}
\def\sS{{\mathbb{S}}}
\def\sT{{\mathbb{T}}}
\def\sU{{\mathbb{U}}}
\def\sV{{\mathbb{V}}}
\def\sW{{\mathbb{W}}}
\def\sX{{\mathbb{X}}}
\def\sY{{\mathbb{Y}}}
\def\sZ{{\mathbb{Z}}}

% Entries of a matrix
\def\emLambda{{\Lambda}}
\def\emA{{A}}
\def\emB{{B}}
\def\emC{{C}}
\def\emD{{D}}
\def\emE{{E}}
\def\emF{{F}}
\def\emG{{G}}
\def\emH{{H}}
\def\emI{{I}}
\def\emJ{{J}}
\def\emK{{K}}
\def\emL{{L}}
\def\emM{{M}}
\def\emN{{N}}
\def\emO{{O}}
\def\emP{{P}}
\def\emQ{{Q}}
\def\emR{{R}}
\def\emS{{S}}
\def\emT{{T}}
\def\emU{{U}}
\def\emV{{V}}
\def\emW{{W}}
\def\emX{{X}}
\def\emY{{Y}}
\def\emZ{{Z}}
\def\emSigma{{\Sigma}}

% entries of a tensor
% Same font as tensor, without \bm wrapper
\newcommand{\etens}[1]{\mathsfit{#1}}
\def\etLambda{{\etens{\Lambda}}}
\def\etA{{\etens{A}}}
\def\etB{{\etens{B}}}
\def\etC{{\etens{C}}}
\def\etD{{\etens{D}}}
\def\etE{{\etens{E}}}
\def\etF{{\etens{F}}}
\def\etG{{\etens{G}}}
\def\etH{{\etens{H}}}
\def\etI{{\etens{I}}}
\def\etJ{{\etens{J}}}
\def\etK{{\etens{K}}}
\def\etL{{\etens{L}}}
\def\etM{{\etens{M}}}
\def\etN{{\etens{N}}}
\def\etO{{\etens{O}}}
\def\etP{{\etens{P}}}
\def\etQ{{\etens{Q}}}
\def\etR{{\etens{R}}}
\def\etS{{\etens{S}}}
\def\etT{{\etens{T}}}
\def\etU{{\etens{U}}}
\def\etV{{\etens{V}}}
\def\etW{{\etens{W}}}
\def\etX{{\etens{X}}}
\def\etY{{\etens{Y}}}
\def\etZ{{\etens{Z}}}

% The true underlying data generating distribution
\newcommand{\pdata}{p_{\rm{data}}}
% The empirical distribution defined by the training set
\newcommand{\ptrain}{\hat{p}_{\rm{data}}}
\newcommand{\Ptrain}{\hat{P}_{\rm{data}}}
% The model distribution
\newcommand{\pmodel}{p_{\rm{model}}}
\newcommand{\Pmodel}{P_{\rm{model}}}
\newcommand{\ptildemodel}{\tilde{p}_{\rm{model}}}
% Stochastic autoencoder distributions
\newcommand{\pencode}{p_{\rm{encoder}}}
\newcommand{\pdecode}{p_{\rm{decoder}}}
\newcommand{\precons}{p_{\rm{reconstruct}}}

\newcommand{\laplace}{\mathrm{Laplace}} % Laplace distribution

\newcommand{\E}{\mathbb{E}}
\newcommand{\Ls}{\mathcal{L}}
\newcommand{\R}{\mathbb{R}}
\newcommand{\emp}{\tilde{p}}
\newcommand{\lr}{\alpha}
\newcommand{\reg}{\lambda}
\newcommand{\rect}{\mathrm{rectifier}}
\newcommand{\softmax}{\mathrm{softmax}}
\newcommand{\sigmoid}{\sigma}
\newcommand{\softplus}{\zeta}
\newcommand{\KL}{D_{\mathrm{KL}}}
\newcommand{\Var}{\mathrm{Var}}
\newcommand{\standarderror}{\mathrm{SE}}
\newcommand{\Cov}{\mathrm{Cov}}
% Wolfram Mathworld says $L^2$ is for function spaces and $\ell^2$ is for vectors
% But then they seem to use $L^2$ for vectors throughout the site, and so does
% wikipedia.
\newcommand{\normlzero}{L^0}
\newcommand{\normlone}{L^1}
\newcommand{\normltwo}{L^2}
\newcommand{\normlp}{L^p}
\newcommand{\normmax}{L^\infty}

\newcommand{\parents}{Pa} % See usage in notation.tex. Chosen to match Daphne's book.

\DeclareMathOperator*{\argmax}{arg\,max}
\DeclareMathOperator*{\argmin}{arg\,min}

\DeclareMathOperator{\sign}{sign}
\DeclareMathOperator{\Tr}{Tr}
\let\ab\allowbreak


\usepackage{hyperref}
\usepackage{url}
\usepackage{graphicx}
\usepackage{float}


% \title{A quantitative analysis of semantic information in deep representations of text and images}
\title{Placeholder for a less complex title}

% Authors must not appear in the submitted version. They should be hidden
% as long as the tmlr package is used without the [accepted] or [preprint] options.
% Non-anonymous submissions will be rejected without review.

\author{\name Santiago Acevedo\footnote{These authors contributed equally to this work.\label{fn:repeat}} \email{sacevedo@sissa.it} \\
      \addr {Scuola Internazionale Superiore di Studi Avanzati (SISSA)}
      \AND
      \name Andrea Mascaretti\footnotemark[\ref{fn:repeat}] \email{amascare@sissa.it} \\
      \addr {Scuola Internazionale Superiore di Studi Avanzati (SISSA)}
      \AND
      \name Riccardo Rende \email{rrende@sissa.it} \\
      \addr {Scuola Internazionale Superiore di Studi Avanzati (SISSA)}
      \AND
      \name Matéo Mahaut \email{mateo.mahaut@upf.edu} \\
      \addr{Universitat Pompeu Fabra (UPF)}
      \AND
      \name Marco Baroni \email{marco.baroni@upf.edu} \\
      \addr{Catalan Institute of Research and Advanced Studies (ICREA)}
      \AND
      \name Alessandro Laio \email{laio@sissa.it} \\
      \addr {Scuola Internazionale Superiore di Studi Avanzati (SISSA)}
}

% The \author macro works with any number of authors. Use \AND 
% to separate the names and addresses of multiple authors.

\newcommand{\fix}{\marginpar{FIX}}
\newcommand{\new}{\marginpar{NEW}}

\def\month{MM}  % Insert correct month for camera-ready version
\def\year{YYYY} % Insert correct year for camera-ready version
\def\openreview{\url{https://openreview.net/forum?id=XXXX}} % Insert correct link to OpenReview for camera-ready version


\begin{document}

\maketitle

\begin{abstract}

\end{abstract}

\section{Introduction}
% Current research explores the conjecture that Large Language Models (LLMs) and foundation vision models such as Vision Transformers (ViTs) possess a \textit{universal language of thought}, enabling them to process concepts across different languages~\citep{peng2024conceptspacealignmentmultilingual, anthropic_claude, platonic_hypothesis, singh-2019-bert, brinkmann2025largelanguagemodelsshare}. A compelling example comes from Anthropic’s study of the Claude model~\citep{anthropic_claude}, which shows that the network occasionally reasons within a shared conceptual space across languages. \citet{platonic_hypothesis} shows that representations of many different LLM/ViTs processing semantically equivalent data are similar to each other, and that this form of alignment is correlated with model size.
% More generally, representations in deep networks, including but not limited to LLMs, are converging, potentially pointing toward shared underlying conceptual structures. 

% In this work, \textit{we seek to turn the qualitative notion that deep networks develop semantically meaningful shared representations into quantitative measures, identifying where in the network, and to what extent, semantic information is encoded}. 
% We can cast the identification of shared semantic content as a problem of relative information: different networks learn  distributions defined on different data domains, such as images and captions. Those distributions loosely span a common  underlying subspace, associated with shared semantic information. The support of the different distributions does not reduce to this subspace, but spans other subspaces, which are data- and architecture-specific. To quantify relative information we need a similarity criterion that is both (i) asymmetric, as there may be a partial order relation between models and representations (qualitatively, the subspaces which are data-specific might have different dimensions); (ii) computationally efficient for representations of dimension of  order $10^5 - 10^6$, which is the typical number of neurons in a deep representation. Cross-entropy is a measure of how difficult it is to encode an event from a distribution into another and it would be the ideal choice, but estimating it is computationally difficult due to the dimension of the representations. On the other hand, local methods like the average number of nearest neighbors \citep{doimo2020,platonic_hypothesis} fail to capture the asymmetry between spaces. To overcome the limitations, we employ the Information Imbalance \citep{II}, a method that leverages conditional ranking to provide an asymmetric measure of relative mutual predictivity, which has been shown to be an excellent proxy of the  cross entropy \citep{Causality-II}.

% We analyze the representation of translations of the same sentences in different languages, using one of the most powerful open-source language models, DeepSeek-V3 with 671B parameters \citep{deepseekv3}, comparing it to the mid-sized Llama3.1-8b LLM \citep{Meta:2024}. Moreover, we analyze the deep representations of pairs of images in ViTs and of the corresponding human-generated descriptions in LLMs, and of pairs of images depicting the same object in ViTs, with the goal of capturing the alignment between image and caption representations in different architectures and  the relative predictive power of images and texts. As visual models, we employ image-gpt-large~\citep{pmlr-v119-chen20s} and DinoV2-large~\citep{oquabDINOv2LearningRobust2024}.

% Our main contributions are as follows:
% \begin{itemize}
%     \item We analyze the DeepSeek-V3 language model and measure, layer by layer, the relative information content between representations of sentence pairs that are translations of each other. This reveals a broad region of the network, robust across language pairs, that consistently encodes shared semantic content. 
%     \item We compare different summaries of the data -considering concatenations, mean and last tokens- to quantify their impact on similarity scores. We find that considering several tokens yields significantly stronger similarity, suggesting that semantic information is spread across many tokens.
%     \item We also identify semantic layers in vision transformers by processing pairs of images that share the same class.
%     These semantic layers are also the most informative about DeepSeek-V3 representations of descriptive image captions, where we observe significant information asymmetries between image and text representations. 
%     \end{itemize}

% \paragraph{Related work.}
% The presence of semantic layers containing shared information in deep networks is connected to the notion of representation alignment~\citep{rep_alignment}. Two different models can be `stitched' together with~\citep{Stitching_bansal} or without~\citep{moschella2023} a trainable stitching layer at several different stages of the network. For instance, two vision models can effectively use similar low-level features~\citep{stitchin0}. Nonetheless, the types of features that we probe for in a semantic layer are instead abstract, encoding the collective relationships between the input constituents, independently of the particular way in which they were expressed in the specific input modality at disposal. 

% These semantic representations lie behind the so-called \textit{Platonic representation hypothesis}, postulating that, as model quality improves, models converge to similar deep representations \citep{platonic_hypothesis}. 

% We significantly extend our understanding of the nature of these shared representations with respect to \citep{platonic_hypothesis}, since for autoregressive language models, the authors take the average across the token axis, reducing the dimensionality of the problem, whereas for vision models, they only consider the ``class'' token. We address instead where semantic information is encoded within the network, namely across different tokens and different layers.

% LLM representations have also been studied in~\citet{cheng2025emergence} using, 
% among other techniques, the Information Imbalance. This work is focused on the analysis of the last-token representation of the same text in different models, and not on multimodal data.

% Other earlier works point to deep similarities across architectures and
% modalities~\citep{Kornblith:etal:2019,Sorscher:etal:2022,Maniparambil:etal:2024}.
% Again, these works do not focus on determining the location and nature
% of the shared semantic layers, nor do they include a latest-generation,
% large model such as DeepSeek-V3. Our work also relates to studies that
% probe the extent to which multilingual LLMs produce shared
% representations across languages (and whether this is due to English
% acting as a pivot language), e.g.\
% \citep{peng2024conceptspacealignmentmultilingual,singh-2019-bert,
% brinkmann2025largelanguagemodelsshare,llama_english,llm_multilingualism}.
% In particular, \citet{llama_english} analyzes whether multilingual
% models align non-English inputs by implicitly routing them through
% English, focusing on logit lens experiments.

% To this research line, we contribute a method to compute similarity that allows us analyze  long spans of tokens and to capture asymmetry between
% language-specific representations, without the use of the embedding or unembedding matrices. Finally, \citet{chughtai2023toy}
% also study universality in learned representations, but while their
% analysis is purely theoretical and operates in controlled toy settings,
% our work examines analogous questions using real-world multilingual and
% multimodal data on actual state-of-the-art LLMs and ViTs.


%\subsection{The Information Imbalance}


Large transformers encode information in high-dimensional spaces,
transforming representations across layers to perform a task. The
Platonic Representation Hypothesis \citep{platonic_hypothesis}
suggests that--at large model sizes--representations of semantically related inputs converge to
similar neighboring structures regardless of (i) the task of the
model and (ii) the specific encoding of the information. The emerging
regularity of these structures underpins the idea that meaning drives
their organization, with models loosely acting as continuous mappings
from a hidden manifold in which semantically similar concepts are close
to each other. We can thus cast the convergence to a universal
representation as a mutual information problem. If we model pairs of representations as random variables with a joint probability distribution, for instance an image and its
caption, we can measure the information lost by replacing the joint
distribution $p(x,y)$ with the product of the marginals $p(x)p(y)$.
If the distributions are independent, observing captions is not
informative of the density of images. If representations converge, a caption will be highly informative about the respective image, and the joint distribution is
not approximated well by the marginals.

Mutual information alone, however, is not sufficient to compare
different models. 
%The Platonic Representation Hypothesis argues that convergence towards a universal set of representations depends on model size. We therefore need to determine (i) whether the effect correlates with the specific encoding of the information (images, captions, \ldots); (ii) whether architectural differences matter; and (iii) how convergence changes with model size. From an information-theoretic standpoint, we need an asymmetric measure that quantifies how much uncertainty over the representation in $Y$ (e.g., captions) remains once we observe the representation in $X$ (e.g., images). 
 Tho distributions $p(X)$ and $p(Y)$ loosely span a common  underlying subspace, associated with shared semantic information. The support of the different distributions does not reduce to this subspace, but spans other subspaces, which are data- and architecture-specific. To quantify relative information we need a similarity criterion that is both (i) asymmetric, as there may be a partial order relation between models and representations (qualitatively, the subspaces, which are data-specific, might have different dimensions); (ii) computationally efficient for representations of dimension of order of thousands, which is the typical number of neurons in a deep representation. Cross-entropy is a measure of how difficult it is to encode an event from a distribution into another and it would be the ideal choice, but estimating it is computationally difficult due to the dimension of the representations. On the other hand, local methods like the average number of nearest neighbors \citep{doimo2020,platonic_hypothesis} fail to capture the asymmetry between spaces. To overcome these limitations, we employ here the Information Imbalance \citep{II} a rank-based measure that tests whether
neighborhoods in one representation space predict neighborhoods in
another: if representations converge, the nearest neighbors of a
point in $Y$ should also be nearest neighbors of the corresponding
point in $X$ in the other architecture.

Our main contributions are as follows:
\begin{itemize}
\item We show that the Information Imbalance,  a statistical measure already successfully used to perform feature selection\cite{} and to detect the presence of causal links\cite{}, can be used to assess quantitatively the information provided by a representation on another representation. We compare its properties to those of other statistics used in the literature for similar purposes. 
\item We focus on DeepSeek-V3 , the largest publicly available LLM, and measure, layer by layer, the relative information content between representations of sentence pairs that are translations of each other. This reveals a region of the network, robust across language pairs, that consistently encodes shared semantic content. We find moreover clear asymmetric effects, such that English-sentence representations tend to be better predictors of their translation than the other way around, and DeepSeek-V3 better predicts smaller Llama3 than the reverse.
\item We compare different aggregation methods for sequence representations, exploring concatenation, mean and last tokens, in order to quantify their impact on similarity scores. We find that utilizing several tokens yields significantly stronger similarity, suggesting that semantic information is spread across many tokens. Furthermore, averaged representations give better predictability scores than concatenated ones, probably due to the ablation of semantically irrelevant positional information.
\item We also identify semantic layers in vision transformers by processing pairs of images that share the same class. These semantic layers are also the most informative about DeepSeek-V3 representations of descriptive image captions, where we observe significant information asymmetries between image and text representations

\end{itemize}

\subsection{Related work}
\textbf{Representation alignment and convergence.} A growing body of work suggests that deep networks develop similar representational structure across architectures and modalities~\citep{Kornblith:etal:2019, Sorscher:etal:2022, Maniparambil:etal:2024, rep_alignment}. This is formalized by the \textit{Platonic representation hypothesis}~\citep{platonic_hypothesis}, which postulates that, as model quality improves, representations converge to a shared structure. Evidence for this includes the finding that models can be stitched together~\citep{stitchin0,Stitching_bansal,moschella2023} across architectures. However, these works largely treat representations as monolithic objects, without asking \emph{where} in the network shared structure emerges, or \emph{how much} information one representation carries about another.

\textbf{Multilingual representations.} 
% A related question is whether multilingual LLMs develop language-independent representations~\citep{peng2024conceptspacealignmentmultilingual, singh-2019-bert, brinkmann2025largelanguagemodelsshare, llm_multilingualism}, and whether this arises from English acting as a pivot language~\citep{llama_english}. 
The related question of whether multilingual LLMs develop language-independent representations was partially studied by searching for a linear transformation between simple concepts expressed in different languages~\citep{peng2024conceptspacealignmentmultilingual} and by using Sparse Auto Encoders~\citep{brinkmann2025largelanguagemodelsshare}. 
%
Another complementary line of research studied the role of English acting as an internal pivot language using logit lens~\citep{llama_english} and neuron importance scores~\citep{llm_multilingualism}, to try to explain the observed representation similarities across languages.
%
\citet{cheng2025emergence} study LLM representations layer by layer using, among other tools, the Information Imbalance, focusing solely on last-token representations of the same text across models. 
% These works, however, do not extend to multimodal data or quantify directional predictability between representations.

\textbf{This work.} 
By measuring the Information Imbalance across different models, layers and token aggregations, we observe that semantic information is spread across many tokens of deep internal representations, and we quantify the \emph{directed} predictability between representations of translations, images sharing a class, and image-caption pairs, unifying and extending recently observed results. 
% We address both gaps. 
% By measuring the Information Imbalance across layers and token aggregations, we identify \emph{which} layers encode shared semantic content and quantify the directed predictability between representations of translations, images sharing a class, and image-caption pairs, without relying on the embedding or unembedding matrices, and at the scale of state-of-the-art models such as DeepSeek-V3.

\section{Methods}
\subsection{Representation choice}
\textbf{Text.} LLMs encode sentences (or other text units) as a sequence of tokens. What we refer to as the semantic content (or meaning) of a sentence should be a global property emerging from the interaction of all the tokens present in it, encoded in the respective representations.  
However, in deep layers, the attention mechanism can move information across tokens, and eventually concentrate it in a specific token, as indeed happens to some extent in the last token of autoregressive LLMs.
%by construction in the last layer of causal models, where the last token representation must encode all information used for next-token prediction.
Thus, it is natural to ask how the choice of token representation aggregation affects the observed representation alignments. 
In Sec.~\ref{sec:last-concat-avg}, we consider using the last token representation as a proxy for the whole sentence, and compare it to two different ways of aggregating token sequences: their average and their concatenation. Given that the averaged representation provided the best alignment scores, we used it in all subsequent experiments. 


\subsection{Data and Models}

\textbf{Translations.} For experiments with translations, we use samples from Opus Books~\citep{opus_books}, with pairs of sentences with lengths between 40 and 80 tokens, extracted from different novels.
We consider between one and three k thousand pairs of sentences where always one of the languages is fixed to English, and the other language is one among Spanish, Italian, German, French, Dutch, or Hungarian, chosen by their presence in the repository with more than a thousand samples after filtering, support by the language models we use, and our ability to manually check translation quality.
%
For text representations, we focus on DeepSeek-V3, a MoE model with 671B total parameters~\citep{deepseekv3}; to compare it with a widely used family of smaller dense models, we also use LLama3~\citep{Meta:2024} with 1, 3 and 8 billion parameters.

\textbf{Images and captions.} 

We consider two architectures with fundamentally different training objectives: image-gpt-large~\citep{pmlr-v119-chen20s} and DinoV2-large~\citep{oquabDINOv2LearningRobust2024}. Image-gpt performs next-token prediction after down-scaling, row-wise unrolling, and color quantization of the input image, closely mimicking the autoregressive strategy of LLMs; previous work suggests that its semantic representations peak in the middle of the network~\citep{pmlr-v119-chen20s}. DinoV2 trains a student network to match a teacher on different augmentations of the same image, producing representations whose final layer is designed as the input to downstream tasks such as depth estimation and image segmentation. 
% We represent both architectures through their mean-token activations.

For experiments with images, we use the ImageNet-1k dataset~\citep{Deng:etal:2009}, sampling 1,000 same-class pairs without replacement (to ensure we have no duplicates in the data) as a proxy for semantically related content, processed with image-gpt-large and DinoV2-large. 
For image--caption experiments, we use the Flickr30k dataset~\citep{youngImageDescriptionsVisual2014}, in which each image is paired with five human-generated captions that we concatenate into a single description. Images are processed with image-gpt-large and DinoV2-large; captions are encoded with DeepSeek-V3. We additionally include the ViT-L/14 visual encoder of CLIP~\citep{radford2021learning} and three DinoV2 variants (small, base, large) to probe the effect of joint text-image contrastive training and vision model scale, respectively.

\begin{figure}[t]
\centering
\includegraphics[width=\linewidth]{figs-tmlr/figure1_ii_cka.pdf}
\caption{
  \textbf{a)}~Information Imbalance (II) $\Delta(X\!\to\!Y)$ and $\Delta(Y\!\to\!X)$,
  CKA, and Neighborhood Overlap~(NO) for a synthetic Gaussian construction
  in which each index~$r$ generates a pair $(X_r, Y_r)$ via
  $Y_r = B_r X_r + \varepsilon$, with $X_r \sim \mathcal{N}(0, I)$,
  $\varepsilon \sim \mathcal{N}(0, \sigma^2 I)$, in $p = 10$ dimensions.
  The matrices $B_r \in \mathbb{R}^{p \times p}$ have monotonically
  increasing rank, from one at $r{=}1$ to full rank at the final index.
  Note that \emph{smaller} II means more predictivity, whereas
  \emph{higher} CKA and NO mean higher similarity.
  The Information Imbalance is the only measure able to reveal that $X$
  is more informative than its noisy transform~$Y$.
  \textbf{b),~c)}~Statistical-power benchmark on a high-dimensional
  Gaussian vector with $p = 10^2$ and $p = 10^4$ components.
  We compute II, CKA, and NO with respect to a smaller vector with only a fraction of the~$p$ features, ranging from small subsets to the full vector.
  In all panels, shaded bands show the standard error over ten
  jackknife repetitions.
  }
\label{fig:mi_ce_vs_ii}
\end{figure}

\section{Results}

\subsection{Quantifying the relative information content between representations}
Our analysis is focused on the quantification of the relative information content of deep representations of translations of the same sentences, images with similar meaning, and image-caption pairs. 
We perform this analysis using the Information Imbalance~\cite{II,Causality-II}, a statistics which (qualitatively) measures the capability of a representation to predict another representation. 
Before discussing our main results, we present II measurements on synthetic data distributions. 
For a comparison, we analyze the same data with other standard statistics used to compare different representations, in particular the Central Kernel Alignment (CKA)~\citep{CKA} and the Neighborhood Overlap~\citep{platonic_hypothesis}.

The Information Imbalance compares the neighborhood
structures of two representation spaces $X$ and $Y$. Given
representations $\{\mathbf{z}_i^X\}_{i=1}^N$ and $\{\mathbf{z}_i^Y\}_{i=1}^N$,
we compute all pairwise distances and rank points $j \neq i$ in each
space by increasing distance, obtaining ranks $r_{i,j}^X$ and
$r_{i,j}^Y$. The Information Imbalance from $X$ to $Y$ is the
normalized average rank in $Y$ of the nearest neighbor of $i$ in $X$:
\[
\Delta(X \rightarrow Y)
= \frac{2}{N-1}\,\frac{1}{N}\sum_{i,j: r^X_{ij=1} } r_{i,j}^{Y},
\]
where the sum is taken over pairs of indices $i,j$ such that $j$ is the first neighbor of $i$ in space $X$, i.e., $r^X_{ij}=1$. 
%where $\mathrm{NN}(i)$ denotes the index of the nearest-neighbor to $i$ in space $X$, i.e., $r_{i,\mathrm{NN}(i)}^X = 1$. 
%
If neighborhoods in
$X$ predict neighborhoods in $Y$, these ranks are small and
$\Delta(X \rightarrow Y)$ is close to zero; if $X$ carries no
information about $Y$, the expected rank is uniform and
$\Delta(X \rightarrow Y) \approx 1$. The asymmetry
$\Delta(X \rightarrow Y) \neq \Delta(Y \rightarrow X)$ quantifies
directional predictability across models or modalities~\citep{II}.



We first compute the Information Imbalance, the CKA and the Neighborhood Overlap for a synthetic Gaussian construction in which each index $r$ generates a pair $(X_r, Y_r)$ via $Y_r = B_r X_r + \varepsilon$, with $X_r \sim \mathcal{N}(0, I)$ and $\varepsilon \sim \mathcal{N}(0, \sigma^2 I)$, in $D = 10$ dimensions. The $B_r$ matrices are defined by first choosing a target rank $r \in \{1,\dots,D\}$ and then sampling independent Gaussian factors $U_r \sim \mathcal{N}(0,1)^{D \times r}$ and $V_r \sim \mathcal{N}(0,1)^{r \times D}$, with the linear map given by $B_r = U_r V_r$, which enforces a transformation of rank $r$. High–rank settings preserve most of the structure in $X_r$, whereas low–rank settings collapse information into a low-dimensional, noisy subspace. 
In Fig.~\ref{fig:mi_ce_vs_ii}-a), we show that the Information
Imbalance captures both the strength and the direction of the relative information
across layers in this controlled synthetic setting. 
In particular, the  Information-Imbalance shows that the predictivity between $X$ and $Y$ is \emph{directional}. 
Indeed, $\Delta(X\!\to\!Y)$ is significantly lower than $\Delta(Y\!\to\!X)$. For comparison, we also report CKA and the Neighborhood overlap, which track the overall loss of alignment between $X_r$ and $Y_r$ but, being symmetric, cannot expose the asymmetry captured by Information Imbalance. 

A complementary benchmark is shown in Figs.~\ref{fig:mi_ce_vs_ii}-b) and~\ref{fig:mi_ce_vs_ii}-c), 
where we compute the Information Imbalance between a Gaussian vector with $p$ components and a smaller dimensional vector containing only a given fraction of them, for cases $p = 10^2$ and $p = 10^4$, respectively. 
For a $10^4$-dimensional Gaussian, CKA saturates close to 1 even when only a small fraction of the components is observed, while the Information Imbalance and the Neighborhood Overlap retain discriminative resolution up to a much larger fraction of shared features. 
This illustrative example echoes the findings of~\citet{platonic_hypothesis}, that observed their Neighborhood Overlap was preferable to CKA in high dimensions.
%
We highlight that in Figs.~\ref{fig:mi_ce_vs_ii}-b) and~\ref{fig:mi_ce_vs_ii}-c) the Information Imbalance is more sensitive in the regime of small fraction of shared features (weak signal) than the Neighborhood Overlap, and complementary, the opposite happens in the regime of large fraction of shared features. In both panels, the information asymmetry between representations is captured by the small gap between Information Imbalance curves.

\subsection{Translations}

We now present the results of the analysis performed using the Information Imbalance on deep representations of data points which, qualitatively, convey approximately the same information to a human. We start by considering  translations of the same sentence in different languages. 

\subsubsection{Representation choice}
\label{sec:last-concat-avg}

Fig.~\ref{fig:last-concat-avg}-a) shows the II between English and Italian translations from the Opus Books parallel corpus collection, using three different representations: (i) The last token, (ii) the concatenation of the last $T$ tokens, or (iii) their average. 
Note that -under the Platonic representation hypothesis- for an ideal dataset of translations, an ideal network separately processing the sentence pairs should generate equivalent representations, namely somewhere in the network the II from one language to the other should be close to zero. Thus, we expect better abstract representations to have a lower Information Imbalance.
In all three cases, the II is minimum in central layers, that were independently found to be related to an abstraction phase\citep{cheng2025emergence}, simultaneously far away from the encoding $(l=0)$ and decoding $(l=L)$ layers, which necessarily contain more language-specific information.
%
We highlight that the predictability between representation significantly increases from single-token to many-token representations. Given that the observed neighborhood predictability seems driven by semantic correspondences, this result suggests that semantic information is spread across many tokens, and not concentrated in, say, only the last token.
%

Among the two considered aggregation methods, i.e., concatenation and average of the last $T$ tokens, we note three main differences. 
In initial layers, the average of an increasing number of tokens increases predictability (reducing II) between representations with respect to the last token alone, whereas this does not happen for the concatenated representation. 
This is consistent with the fact that the meaning of sentences, which drives predictability of the representations, does not rely on precise positional information and it will be typically distributed across different tokens in a sentence and its translation. 
Furthermore, we observe that, in deep layers, the average of a larger number of tokens in the representation monotonically improves predictability (whereas this does not happen for the concatenated representations) and, most importantly, gives the ``best'' (i.e., lowest) predictability score. 
Consequently, in what follows we focus on token-averaged representations.

\begin{figure}[t]
    \centering
    \includegraphics[width=\linewidth]{figs-tmlr/translations/deepseek-ii-english-italian-avg-vs-concat-panels.pdf}
    \caption{\textbf{Representation choice.}
    Information Imbalance from English to Italian tranlsations in DeepSeek-V3 when tokens are a) averaged or b) concatenated. The (hardly visible) shaded colored areas correspond to the standard deviation obtained by subsampling half of the samples five times. }
    \label{fig:last-concat-avg}
\end{figure}

Note that, for all curves, even if predictability is highest in the central layers, it is still larger at the last layer than at the first. This is probably due to the fact that, at the last layer, the representation must still contain some of the semantic information that was computed in the central layers.
%
As a null hypothesis, we show in Supp. Inf.~\ref{sec:non-informative-check} that performing batch-shuffling on any of the datasets leads to completely uninformative representations (Information Imbalance givess 1 at every layer), since the semantic correspondence between translations is destroyed. 


\begin{figure}
    \centering
    \includegraphics[width=\linewidth]{figs-tmlr/translations/concat-english-pivot-llamas.pdf}
    \caption{\textbf{Comparison with other datasets and models.}
    Panel a): Information Imbalance from English to several languages, computed on representations generated by DeepSeek-V3.
    Panel b): Information Imbalance from English to Italian, computed on representations generated by Llama3 models with 1,3, and 8 billion parameters, and by DeepSeek-V3 for comparison. 
    In both panels we used the average of the last 20 tokens, and the (hardly visible) shaded areas correspond to the standard deviation obtained subsampling half of the samples five times. 
    }
    \label{fig:compare-langs-models}
\end{figure}

\subsubsection{Comparison with other datasets and models}

Fig.~\ref{fig:compare-langs-models}-a) shows that the II profiles across layers for Spanish, Dutch, German and French datasets of translations are in very good agreement with the Italian one. Hungarian, which is simultaneously the tested language with the smallest amount of resources available on the internet\footnote{\url{https://en.wikipedia.org/wiki/Languages_used_on_the_Internet}} and the only non-Indo-European language in the set, shares the same qualitative II profile, but with  larger II values (smaller predictability). 
%
Fig.~\ref{fig:compare-langs-models}-b) shows instead that the II layer profile for English-Italian translations is similar for the Llama3 family of models with 1, 3 and 8 billion parameters. DeepSeek-V3 results from Fig.~\ref{fig:last-concat-avg}-a) are added for direct comparison. Notably, and in accordance with the Platonic representation hypothesis, increasing model size reduces the II, predominantly in the inner layers. The position of the global II minimum also shifts slightly towards earlier layers for the larger models. We highlight that, qualitatively, the II profiles of all models are very similar, and DeepSeek-V3's global minimum is deeper than for the other models. 

%In  Supp.~Inf.~\ref{sec:comparison_alignment} we show that our results are compatible with those obtained using the overlap neighborhood metric from~\citet{platonic_hypothesis}.


\subsubsection{Information asymmetries}
\label{sec:inf-asym}

Given that English is the dominant language in terms of training resources and performance \citep{Zhu:etal:2024}, one could expect English representations to be more informative than those generated processing other languages.
%
Fig.~\ref{fig:english-pivoting}-a) shows that this asymmetry is indeed observed, most notably in the last quarter of the network, for the specific case of English-Italian translations. 
Nonetheless, the II asymmetry profile is far from trivial. Notably, both representations are equally predictive of each other where predictability is maximal, around the middle of the network. This is in line with the hypothesis that the middle layers store semantic information that might be largely language-independent.  
Moving towards earlier layers, at around $0.2$ relative depth, we see that a small gap opens between the two II curves, where again English is the most informative, whereas in the very first layers we observe a small inverted asymmetry.
Fig.~\ref{fig:english-pivoting}-b) shows the II asymmetry $\mathcal{A} = II(English \rightarrow other) - II(other \rightarrow English)$ across the layers of DeepSeek-V3. Note that with this definition, a negative asymmetry corresponds to representations in English being more informative than representations in the other language. We observe that asymmetry profile across layers is robust for different datasets of translations in five different languages, with the same overall trend (the representations developed from English sentences are asymmetrically more informative in the early and late layers).
%
Fig.~\ref{fig:llama's-A} in Appendix~\ref{sec:llama's-A} shows that the corresponding asymmetry profiles across the layers of Llama3-8b also are positive in the first layers, zero in middle layers and negative in the last quarter of the network, with the exception of the last layer, in which the asymmetry is roughly zero.
%

One can also wonder if the representation of a huge modern model like DeepSeek-V3 is more informative than that generated by a smaller Llama3 model. 
%
For identical English inputs, Fig.~\ref{fig:english-pivoting}-c) shows the II between DeepSeek-V3 and Llama3-8b representations  at comparable relative depths.  In contrast to the II profile between translations in Fig.~\ref{fig:last-concat-avg}, in Fig.~\ref{fig:english-pivoting}-c) the II in the very first layers is very small, suggesting that the word embeddings of the two models are very similar. 
The II has then a small peak at roughly $0.1$ relative depth, plausibly related to some difference in  input processing between the two networks. Afterwards, at intermediate depths, we observe a region of maximal mutual predictability (lowest II) without information asymmetries up until roughly  half of the networks. Again, this is plausibly related to the fact that the middle layers are those that contain the most ``universal'' semantic information. In the second part of the network, we observe a  massive gap between the curves, where DeepSeek-V3 representations are far more predictive than those from Llama3-8b. 

\begin{figure}[t]
    \centering
    \includegraphics[width=\linewidth]{figs-tmlr/translations/English-pivoting-3panel.pdf}
    \caption{\textbf{Information asymmetries.}
    Panel a): Information Imbalance from English to Italian and from Italian to English, computed on representations generated by DeepSeek-V3. 
    Panel b): Asymmetry $\mathcal{A} = II(English \rightarrow other) - II(other \rightarrow English)$ in Information Imbalance between English and other languages, computed on representations generated by DeepSeek-V3. Note that with this definition, negative asymmetry means that English is more informative than the other language.
    Panel c): Information Imbalance between representations generated by DeepSeek-V3 and Llama3-8b on the same English text, as a function of comparable relative depth in both models. 
    The (hardly visible) shaded area corresponds to one standard deviation, computed with a Jackknife procedure subsampling five times half of the samples.
    In all three panels we used the average of the last 20 tokens.
    }
    \label{fig:english-pivoting}
\end{figure}

\subsubsection{Token-token Information Imbalance}
\label{sec:token-token-II}

We have used representations constructed out of the average of several tokens to measure mutual predictability between translations (Fig.~\ref{fig:last-concat-avg}) and to systematically quantify information asymmetries between English and other languages (Fig.~\ref{fig:english-pivoting}-b)).  In this section, we zoom-in on each individual language and look at the correlation structure between tokens, trying to gain further qualitative insights.

Fig.~\ref{fig:corr_ds} shows the II between the last token and a previous token at distance $\tau$ from it, for English and Italian representations generated by both DeepSeek-V3 and Llama3-8b. 
%
For DeepSeek-V3, in the very first layers, the last token is, on average, not informative about a previous token at distance $\tau \approx 5$. 
%
In deeper layers, tokens become more and more informative about each other.  
The layers with strongest correlations (lowest II) are those with relative depth roughly between $0.6$ and $0.9$, which approximately coincides with the broad region with very low II between translations in Fig.~\ref{fig:last-concat-avg}, right after the global minimum.
%
The token-token predictability in these deep layers is much larger in DeepSeek-V3 (Figs.~\ref{fig:corr_ds}-a) and~\ref{fig:corr_ds}-b)) than in Llama3-8b (Figs.~\ref{fig:corr_ds}-c) and~\ref{fig:corr_ds}-d)). 
%
This analysis is to be compared with the one of Fig.~\ref{fig:compare-langs-models}-b), where DeepSeek-V3's II between translations is lower than the II computed on any of the Llama3 models. 
Thus, we observe that representations that are more predictive across languages are also those in which the information is more consistent across many tokens. 

Furthermore, by comparing Fig.~\ref{fig:corr_ds}-a) and Fig.~\ref{fig:corr_ds}-b), we see that  the token-token II in very deep layers of DeepSeek-V3 is also much lower in English than in Italian, which could in turn be connected with the information asymmetries observed between these two languages in Fig.~\ref{fig:english-pivoting}-a).


\begin{figure}[t]
    \centering
    \includegraphics[width=\linewidth]{figs-tmlr/translations/4-panel-corrs_single-row.pdf}
    \caption{\textbf{Token-token Information Imbalance.}
    Information Imbalance from the last token to a previous token at distance $\tau$, as a function of $\tau$. Panels a) and b) correspond to DeepSeek-V3 representations of English and Italian text, whereas panels c) and d) correspond to Llama3-8b representations of English and Italian. 
    The (hardly visible) shaded area corresponds to one standard deviation, computed with a Jackknife procedure subsampling five times half of the samples.
    }
    \label{fig:corr_ds}
\end{figure}


\subsection{Images with related content}

\paragraph{Within-model analysis.}
We first compare representations for pairs of images that share the same ImageNet-1k class but depict \emph{different} scenes, using mean-token activations.
Since the assignment of the two images to the two representation spaces is arbitrary, we report the average Information Imbalance in both directions.
Results are shown in Fig.~\ref{fig:image-image}\,a).
For DinoV2, the Information Imbalance decreases monotonically and reaches values close to~$0.35$ around its last layer.
For ImageGPT, the minimum ($\approx 0.65$) is attained around layer~20 ($42\%$ relative depth, out of 48~layers), after which the Information Imbalance rises again, reaching $\approx 0.85$ at the last layer.
This is consistent with the respective training objectives: DinoV2's self-supervised objective produces increasingly semantic representations toward the output, whereas ImageGPT, like autoregressive LLMs, concentrates semantic content in the middle of the network and reverts to next-token-specific features near the output.

\paragraph{Cross-model analysis.}
We next ask whether the semantic layers identified above are also the most mutually predictive \emph{across} the two architectures.
In Fig.~\ref{fig:image-image}\,b), we fix one model at its most semantic layer---DinoV2 at its last layer, ImageGPT at layer~20---and sweep across all layers of the other model, using the same ImageNet-1k same-class pairs (i.e.\ different images of the same class processed by different models).
The cross-model Information Imbalance confirms that the semantic layers of one architecture are best predicted by the semantic layers of the other: ImageGPT layers around $40\%$ relative depth are most predictable from DinoV2's last layer, and DinoV2's deepest layers are most predictable from ImageGPT layer~20.

Fig.~\ref{fig:image-image}\,c) repeats the cross-model analysis on Flickr30k, where both models encode the \emph{same} images.
The qualitative picture is unchanged: the same semantic regions emerge, and the Information Imbalance profiles are consistent with panels~a) and~b), confirming that the identified semantic layers generalize across datasets.
In both panels~b) and~c), the dashed lines (backward direction) closely track the solid lines (forward), indicating approximately symmetric mutual predictability between the two models at their respective semantic layers.

\begin{figure}[t!]     
      \centering
      \includegraphics[width=\textwidth]{figs-tmlr/image-captions/fig4a.pdf}
      \caption{
      Information Imbalance between ImageGPT and DinoV2 image
      representations, using mean-token activations.
      \textbf{a)}~Within-model Information Imbalance for 1,000 same-class \emph{different}
      image pairs from ImageNet-1k: each model encodes two distinct
      images of the same class, and the Information Imbalance is computed between the two
      resulting representations. DinoV2 is evaluated at its last layer;
      ImageGPT at layer~20.
      \textbf{b)}~Cross-model Information Imbalance on the same ImageNet-1k pairs.
      DinoV2 (last layer) is compared against all ImageGPT layers,
      and ImageGPT (layer~20) against all DinoV2 layers.
      \textbf{c)}~Cross-model Information Imbalance on Flickr30k, where the \emph{same}
      images are encoded by both models. DinoV2 (last layer)
      is swept against all ImageGPT layers, and ImageGPT (layer~20)
      against all DinoV2 layers.
      In~(b) and~(c), solid lines denote the forward direction
      (fixed model$\,\to\,$swept model); dashed lines the backward
      direction (swept model$\,\to\,$fixed model).}
      \label{fig:image-image}
  \end{figure}





\subsection{Images and captions: multimodal data sharing semantic content}

Image--caption pairs convey approximately the same meaning in different modalities, similarly to translations of the same sentence in different languages.
In Fig.~\ref{fig:ds-52}\,a), we report the Information Imbalance between layer~52 of DeepSeek-V3---which lies within the region of maximal mutual predictability between translations (see Fig.~\ref{fig:last-concat-avg})---and each layer of the two vision transformers, on the Flickr30k dataset.
The cross-modal profile recovers exactly the same semantic regions identified by the image--image analysis: the Information Imbalance decreases monotonically toward the last layer for DinoV2, while it reaches a minimum in the middle for ImageGPT.

For DinoV2, the minimum Information Imbalance is $\approx 0.10$ (at layer~24, the final layer), with a reciprocal value of $\approx 0.12$.
The two directions are close, indicating approximately balanced cross-modal predictability; the slight advantage of DeepSeek-V3$\to$DinoV2 over DinoV2$\to$DeepSeek-V3 is consistent with DinoV2 being explicitly trained to extract semantically meaningful visual features, so that its final representation is well predicted by a textual semantic summary.
For ImageGPT, the minimum Information Imbalance from DeepSeek-V3 is $\approx 0.41$ (at layer~20, $42\%$ relative depth), with a reciprocal of $\approx 0.44$.
The substantially higher values compared to DinoV2 indicate weaker cross-modal alignment.
The asymmetry between the two directions, with text predicting image representations slightly better than the reverse, is consistent with DeepSeek-V3 producing higher-quality semantic representations than ImageGPT.

To test whether joint text--image training facilitates cross-modal alignment, we compare in Fig.~\ref{fig:ds-52}\,b) the ViT-L/14 visual encoder of the multimodal CLIP system~\citep{radford2021learning} against both DinoV2-large (a within-modality comparison) and CLIP-Text (a cross-modal comparison).
The minimum Information Imbalance from CLIP-Image to CLIP-Text is $\approx 0.22$ (at layer~19), compared to $\approx 0.05$ from CLIP-Image to DinoV2, showing that two vision models trained with different objectives share more representational structure than a vision and a text encoder, even when jointly trained.

Finally, to probe whether the observed cross-modal alignment depends on the scale of the vision model, we repeat the DeepSeek-V3--DinoV2 analysis with three DinoV2 variants: large (24~layers), base (12~layers), and small (12~layers); see Fig.~\ref{fig:llama-52}\,b).
The minimum Information Imbalance is $0.10$ for DinoV2-large, $0.10$ for DinoV2-base, and $0.15$ for DinoV2-small.
All three profiles reach their minimum at the last layer, confirming that the semantic region is robustly identified regardless of vision model size.
The small gap between the large and base variants, together with the higher value for the small model, suggests that the quality of cross-modal alignment saturates once the vision model reaches a sufficient capacity, consistent with the results on translations in Fig.~\ref{fig:compare-langs-models}.
We report the analysis with LLama3.1-8B as the text encoder in Supp.~Inf.~\ref{appendix:image-captions}.

\begin{figure}[t!]
      \centering
      \includegraphics[width=\textwidth]{figs-tmlr/image-captions/fig4b.pdf}
      \caption{
      Information Imbalance between image and caption representations
      on Flickr30k, using centered and normalized mean-token activations.
      \textbf{a)}~Images encoded with DinoV2-large and ImageGPT; captions
      with DeepSeek-V3 (layer~52). Results shown as a function of
      relative ViT depth.
      \textbf{b)}~CLIP-Image (layer~24, last) compared against DinoV2-large
      (within modality) and CLIP-Text (cross-modal).
      In both panels, solid lines denote the forward direction
      (target$\,\to\,$sweep model); dashed lines denote the backward
      direction (sweep model$\,\to\,$target).
      Uncertainty estimated via bootstrapping
      (100 replicas of 750 samples in~(a);
      2,000 replicas of 200 samples in~(b);
      both over 5,000 total).
      }
      \label{fig:ds-52}
  \end{figure}


\section{Limitations}
\label{sec:limitations}

\section{Conclusion}
\label{sec:conclusions}

What to highlight:

\begin{itemize}
    \item Huth et al.~came up with a sexy hypothesis, but now we need to do some systematic work to systematize it. Sure, there's convergence, but under which conditions, where in the networks, and to what extent it means that two representations become increasingly identical, vs.~one subsumes the other?
    \item In this paper, we started this systematic analysis. In particular, we showed with theoretical experiments that II is an appropriate measure for this type of comparison, and then, for the first time, applied it thoroughly to LLMs and vision models.
    \item First of all, we confirm convergence, and in particular that convergence happens in middle layers of LLMs, whereas for vision transformers it might be more dependent on the training objective.
    \item Cross-linguistic convergence of translations for a large model like DeepSeek-V3 is extremely strong, suggesting genuinely abstract semantic representations in its intermediate layers. Convergence is moreover clearest when sentences are represented through their averaged tokens (suggesting semantic representations are not position-dependent) and taking more tokens into account (suggesting meaning is well-spread across long token spans). At the same time, the better a representational system (better model, more resource-rich language), the larger its similarity across tokens in a sentence, suggesting semantic coherence as a hallmark of representational quality. 
    \item There are clear asymmetries in the expected directions: bigger to smaller model, English to other languages, DeepSeek to vision models. However, in the area of stronger convergence, asymmetry typically goes away or greatly diminishes, suggesting that for all inputs and models we considered there are areas of abstract semantic representations fully comparable between the ``better'' and ``worse'' models and inputs; what changes is how these representations are reached and built upon across layers.
\end{itemize}

\bibliography{marco,refs}
\bibliographystyle{tmlr}

\appendix
\section{Appendix}

\section{Misalignment of translations erases semantic similarity}
\label{sec:non-informative-check}
As a consistency check, Fig~\ref{fig:misalignment} shows the Information Imbalance for DeepSeek-V3's and Llama3.1-8b representations using misaligned translations, namely performing a batch-shuffle in one of the datasets. Since the semantic correspondence between sentences is destroyed, the representations are not informative about each other, and thus the Information Imbalance is close to one, for all layers. The same occurs with misaligned image pairs, Information Imbalance being around one when pairs do not share the same class, for all layers.

\begin{figure}[h!]
    \centering
    \includegraphics[width=0.5\linewidth]{figs.bak/appendix/misalignment_20.pdf}
    \caption{Information Imbalance between English (en) and Spanish (es) representations generated by DeepSeek-V3 and Llama3.1-8b, for the "non-informative case", namely a misaligned dataset in which we batch-shuffle the Spanish translations. 
    The (hardly visible) shaded area corresponds to one standard deviation, computed with a Jackknife procedure subsampling five times half of the samples.}
    \label{fig:misalignment}
\end{figure}

\section{Information Asymmetry in translations processed by Llama3}
\label{sec:llama's-A}
Fig.~\ref{fig:llama's-A} shows the asymmetry profiles across layers, where we observe that 

\begin{figure}[h!]
    \centering
    \includegraphics[width=0.5\linewidth]{figs-tmlr/appendix/Asymmetry-EN-multi-llama.pdf}
    \caption{Caption}
    \label{fig:llama's-A}
\end{figure}

\section{CKA and neighborhood overlap comparison on translations}
\label{CKA-NO-translations}

Fig.~\ref{fig:CKA-translations} compares our II curves between English and Italian translations as a function of DeepSeek-V3's relative depth with CKA and Neighborhood Overlap computed on the same data, using a) the Last token, b) the Average of the last 20 tokens or c), their concatenation.

First, we highlight that in all three representational choices, the three metrics concurrently achieve their maximum score (lowest II, maximum CKA and NO, respectively). 
However, there are several nuances to note. 
%
Starting with the Neighborhood Overlap, we note that, being defined between 0 and 1, here it only reaches maximum values of roughly $0.1$, $0.4$ and $0.2$ for panels a), b) and c), respectively. 
Moreover, the last-layer Neighborhood Overlap is bearly above zero for last token and concatenated token representations, giving $0.031 +- 0.002$ and $0.035 +- 0.003$, respectively.
Small alignment values were also observed and even highlighted as a limitation of the convergence of representations in~\citet{platonic_hypothesis}, although the authors do not provide arguments to understand the reasons behind this effect. 
We note that if the $k$-th neighbor of representation $A$ is instead the $k+1$ neighbor of representation $B$, automatically that point counts as \emph{outside} the $k$-neighborhood, even if it is extremely close, rendering the alignment value very low. 
This effect does not take place in the Information Imbalance, which, instead of evaluating if $k$ neighbors are shared, measures what is the rank in space $B$ of $k$ neighbors in space $A$, and vice-versa.

CKA displays much larger alignment values than the neighborhood overlap in all three panels, with the exception of the last layer of panel a), where CKA = $0.059 +- 0.001$, where it struggles to capture the very weak signal. Moreover, while the CKA similarity profiles are compatible with those computed with II and NO in panels a) and b), we observe that CKA gives very high similarity in initial layers of panel c), in complete contrast with the other two metrics. We hypothesize this to be a spurious artifact of working in extremely high dimensions (20 concatenated tokens, each embedded in 7168 dimensions of DeepSeek-V3), that is avoided by both metrics based on neighborhood ranks.



\begin{figure}[h!]
    \centering
    \includegraphics[width=\linewidth]{figs-tmlr/appendix/CKA-NO-translations.pdf}
    \caption{\textbf{CKA and Neighborhood overlap comparison on translations.}
    Information Imbalance from English to Italian and from Italian to English, together with the CKA and Neighborhood Overlap computed on the same activations, using a) the last token, b) the average of the last 20 tokens and c) the concatenation of the last 20 tokens.
    }
    \label{fig:CKA-translations}
\end{figure}

\subsection{Image-Captions}
\label{appendix:image-captions}
 \begin{figure}[h!]
      \centering
      \includegraphics[width=\textwidth]{figs-tmlr/image-captions/fig4c.pdf}
      \caption{
      Information Imbalance between image and caption representations
      on Flickr30k.
      \textbf{a)}~Images encoded with DinoV2-large and ImageGPT; captions
      with LLama3.1-8B (layer~20). Results shown as a function of
      relative ViT depth.
      \textbf{b)}~Captions encoded with DeepSeek-V3 (layer~52); images
      encoded with DinoV2-large, DinoV2-base, and DinoV2-small.
      In both panels, solid lines denote the forward direction
      (target$\,\to\,$sweep model); dashed lines denote the backward
      direction (sweep model$\,\to\,$target).
      Uncertainty estimated via bootstrapping
      (100 replicas of 750 samples, over 5,000 total).
      }
      \label{fig:llama-52}
  \end{figure}

\end{document}